\chapter{App I --- Feasibility: Task Lists and Risk Assessment}

\section{App I --- Feasibility: Task Lists and Risk Assessment}

This appendix complements Chapter with the full per-phase task lists (task IDs, estimated hours, critical path flags) and the five-phase NIST SP 800-30 risk assessment that underpins the risk matrix in that chapter.

\subsection{Phase Task Lists}

\subsubsection{Phase 0: Preparation (July--November 2025)}

\begin{itemize}
\tightlist
\item
  Task 0.1: HomeLab Initialisation Setup (64h)
\item
  Task 0.2: VSCode Setup for thesis development (18h)
\item
  Task 0.3: Parrot OS Security Setup and Basic Tools (15h)
\item
  Task 0.4: Home Lab As-Built Documentation (8h)
\end{itemize}

\textbf{Duration}: 105 hours

\subsubsection{Phase 1: Preliminary Project (October 2025--January 2026)}

\begin{itemize}
\tightlist
\item
  Task 1.0: Thesis Development (22h)
\item
  Task 1.1: SMART Objectives Definition (4h)
\item
  Task 1.2: State of the Art: Frameworks and Tools (15h)
\item
  Task 1.3: Functional and Non-Functional Requirements (10h)
\item
  Task 1.4: Pentesting Methodology and Validation (8h) --- \textbf{CRITICAL}
\item
  Task 1.5: Design Decisions and Justification (5h)
\item
  Task 1.6: Gantt Chart Development (20h)
\item
  Task 1.7: Risk Analysis and Mitigation Plan (15h)
\item
  Task 1.8: Resource Inventory (5h)
\item
  \textbf{Milestone 1.10}: Preliminary Project (40h) --- \textbf{Delivery: January 16, 2026}
\end{itemize}

\textbf{Duration}: 82 hours

\subsubsection{Phase 2: Interim Report (January--March 2026)}

\begin{itemize}
\tightlist
\item
  Task 2.1: Lab 1 Research --- Reconnaissance (14h) --- \textbf{CRITICAL}
\item
  Task 2.2: Lab 2 Research --- Web Vulnerabilities (14h) --- \textbf{CRITICAL}
\item
  Task 2.3: Lab 3 Research --- Network Analysis/IR (14h) --- \textbf{CRITICAL}
\item
  Task 2.4: Technical Write-ups Labs 1 \& 2 (13h)
\item
  Task 2.5: Topology Adjustments (30h)
\item
  Task 2.6: Pre-validation with Pilot Users (20h) --- \textbf{CRITICAL}
\item
  Task 2.6.1: Pre-validation Revision (6h) --- \textbf{CRITICAL}
\item
  Task 2.7: Analyse and Write Hackathon Results (17h)
\item
  Task 2.8: Interim Report Formalisation (40h)
\item
  Task 2.9: Set Up Servers, VMs, and EVE-NG (20h)
\item
  Task 2.10: Backup and Documentation of EVE-NG (10h)
\item
  \textbf{Milestone 2.11}: Intermediate Thesis (50h) --- \textbf{Delivery: March 20, 2026}
\end{itemize}

\textbf{Duration}: 212 hours

\subsubsection{Phase 3: Final Delivery (March--May 2026)}

\begin{itemize}
\tightlist
\item
  Task 3.1: Lab 4 Research --- Post-Exploitation (16h) --- \textbf{CRITICAL}
\item
  Task 3.2: Lab 3 Validation with Users (12h)
\item
  Task 3.3: Technical Write-ups Labs 3 \& 4 (19h)
\item
  Task 3.4: Penetration Testing Round 1 (Labs 1--2--3) (30h) --- \textbf{CRITICAL}
\item
  Task 3.5: Penetration Testing Round 2 (All Labs) (25h) --- \textbf{CRITICAL}
\item
  Task 3.6: Python Automation Scripts (30h)
\item
  Task 3.7: Final Validation with Pilot Users (18h)
\item
  Task 3.8: Final Review and Project Polishing (25h)
\item
  Task 3.9: Bachelor's Thesis Report Writing I (15h) --- \textbf{CRITICAL}
\item
  Task 3.10: Bachelor's Thesis Report Writing II (40h) --- \textbf{CRITICAL}
\item
  Task 3.11: Bachelor's Thesis Report Writing III (35h) --- \textbf{CRITICAL}
\item
  Task 3.12: Final Revision (15h)
\item
  Task 3.13: Final Report Preparation (20h)
\item
  \textbf{Milestone 3.14}: Final Report (5h) --- \textbf{Delivery: May 27, 2026}
\end{itemize}

\textbf{Duration}: 315 hours

\subsection{NIST SP 800-30 Five-Phase Risk Assessment}

This section details the systematic risk assessment following NIST SP 800-30. The summary risk matrix is provided in Table in Chapter .

\subsubsection{Phase 1: Asset Identification and Scope Definition}

Critical assets are categorised by CIA triad:

\textbf{Hardware Assets}: - HomeLab server (2× Xeon E5-2680v4, 64 GB RAM) --- Availability critical - Development workstations (laptop + desktop) --- Integrity critical - Network equipment (router, switch) --- Availability critical - External storage (backup) --- Confidentiality + Availability critical

\textbf{Software Assets}: - EVE-NG Community Edition --- Availability critical - Parrot OS Security distribution --- Integrity critical - Vulnerable VMs (Metasploitable, DVWA) --- Integrity critical - Development tools (VSCode, Git, LaTeX) --- Integrity critical

\textbf{Data Assets}: - Thesis source code and documentation --- Confidentiality + Integrity critical - Lab topology configurations --- Integrity + Availability critical - Student learning materials --- Availability critical - Project research findings --- Confidentiality + Integrity critical

\textbf{Scope}: Infrastructure and technical risks within the controlled laboratory environment. Excluded: organisational, regulatory, and external threats beyond the academic context.

\subsubsection{Phase 2: Threat and Vulnerability Identification}

\textbf{Threat categories (NIST SP 800-30)}: - \textbf{Hardware Threats}: Disk failures, power supply degradation, CPU overheating, memory exhaustion - \textbf{Software/Configuration Threats}: Outdated VM images, incompatible versions, configuration drift, EVE-NG crashes - \textbf{Knowledge Gap Threats}: Insufficient expertise in advanced networking, virtualisation optimisation, security best practices - \textbf{Resource Constraint Threats}: Limited RAM for simultaneous VM operations, storage capacity limits - \textbf{Integration Threats}: EVE-NG network compatibility issues, tool interoperability - \textbf{Data Loss Threats}: Accidental deletion, corruption, hardware failure

\textbf{Key vulnerabilities}: - Aging server components (end-of-life approaching) - Single point of failure --- no infrastructure redundancy - EVE-NG Community Edition has no commercial support - Limited RAM capacity for production-scale virtualisation - Inadequate documentation of network topology state - Limited backup redundancy (single backup location)

\subsubsection{Phase 3: Impact and Likelihood Analysis}

Scales follow NIST SP 800-30: Probability (Low \textless25 Impact (Low: minor delays, Medium: project delays, High: project jeopardy).

{\footnotesize\begin{longtable}[]{@{}
  >{\raggedright\arraybackslash}p{0.17\linewidth}
  >{\raggedright\arraybackslash}p{0.17\linewidth}
  >{\raggedright\arraybackslash}p{0.17\linewidth}
  >{\raggedright\arraybackslash}p{0.17\linewidth}
  >{\raggedright\arraybackslash}p{0.17\linewidth}@{}}
\toprule\noalign{}
\begin{minipage}[b]{\linewidth}\raggedright
\textbf{Risk Scenario}
\end{minipage} & \begin{minipage}[b]{\linewidth}\raggedright
\textbf{Prob.}
\end{minipage} & \begin{minipage}[b]{\linewidth}\raggedright
\textbf{Impact}
\end{minipage} & \begin{minipage}[b]{\linewidth}\raggedright
\textbf{Level}
\end{minipage} & \begin{minipage}[b]{\linewidth}\raggedright
\textbf{Primary Asset}
\end{minipage} \\
\midrule\noalign{}
\endhead
\bottomrule\noalign{}
\endlastfoot
EVE-NG Performance Issues & Medium & High & HIGH & Infrastructure, Labs \\
VM Image Incompatibility & Medium & Medium & MEDIUM & Development, Labs \\
Knowledge Gaps (advanced net.) & Medium & Medium & MEDIUM & Development, Timeline \\
Hardware Degradation / Disk Failure & Low & High & MEDIUM & Infrastructure \\
Data Loss (insufficient backup) & Low & High & MEDIUM & Documentation \\
Configuration Documentation Gaps & Medium & Medium & MEDIUM & Operations \\
\end{longtable}}

\subsubsection{Phase 4: Risk Evaluation and Prioritisation}

\textbf{HIGH Priority}: EVE-NG Performance Issues --- infrastructure bottlenecks causing lab functionality failures. Medium probability, high impact. Requires immediate resource optimisation and monitoring implementation.

\textbf{MEDIUM Priority}: - VM Incompatibility --- disparate image formats causing integration problems - Knowledge Gaps --- experience gaps in advanced networking/virtualisation - Configuration Documentation --- inadequate topology documentation - Data Loss --- insufficient backup redundancy (single point of failure)

\textbf{LOW Priority}: Hardware Degradation --- aging components manageable through preventive maintenance.

\subsubsection{Phase 5: Mitigation Strategies}

\textbf{EVE-NG Performance Issues (REDUCTION)}: - Hardware dimensioning: 64 GB RAM dedicated to the server running EVE-NG (32 GB proved insufficient --- additional module acquired, 80\,€) - Resource optimisation: memory tuning, VM suspension when idle - Documented fallback: VirtualBox as alternative hypervisor if critical failure - \emph{Result}: Probability reduced from Medium to Low

\textbf{VM Incompatibility (PREVENTION)}: - Early testing: validate all images in target environment before use - Format standardisation: QCOW2 for EVE-NG compatibility - Alternative repositories maintained (VulnHub, GitHub) - \emph{Result}: Probability reduced from Medium to Low

\textbf{Knowledge Gaps (PREVENTION + REDUCTION)}: - Structured learning: Task 1.2 (15h) comprehensive framework research - Continuous upskilling during lab research phases (Tasks 2.1--3.1) - Expert consultation with tutor Pere Vidiella i Catalan - Technical references: NIST, OWASP - \emph{Result}: Impact reduced from Medium to Low-Medium

\textbf{Configuration Documentation (PREVENTION)}: - Continuous documentation: update topology diagrams during development - Automated export: EVE-NG built-in topology export - Version control: all network configs via Git repository - \emph{Result}: Probability reduced from Medium to Low

\textbf{Data Loss (REDUCTION)}: - Daily automated encrypted backups to external storage - Version control: full Git history on GitHub - Off-site cloud backup (student account) - \emph{Result}: Probability Low, Impact Low

\textbf{Hardware Degradation (PREVENTION)}: - Quarterly hardware diagnostics (disk health, temperature monitoring) - Automated S.M.A.R.T. alerts for disk errors and CPU temperature thresholds - Component procurement plan with 6-month horizon for critical parts

\textbf{Risk Assessment Conclusion}: All identified risks have been systematically evaluated and assigned concrete mitigation strategies. No individual risk is considered project-blocking.

\subsection{Detailed Budget Tables}

\subsubsection{Software and Licences}

{\footnotesize\begin{longtable}[]{@{}>{\raggedright\arraybackslash}p{0.22\linewidth}
  >{\raggedright\arraybackslash}p{0.22\linewidth}
  >{\raggedright\arraybackslash}p{0.22\linewidth}
  >{\raggedright\arraybackslash}p{0.22\linewidth}@{}}
\toprule\noalign{}
\textbf{Resource} & \textbf{Licence Type} & \textbf{Unit} & \textbf{Total} \\
\midrule\noalign{}
\endhead
\bottomrule\noalign{}
\endlastfoot
EVE-NG Community Edition & Free/Open Source & 0\,€ & 0\,€ \\
Parrot OS Security & Free/Open Source & 0\,€ & 0\,€ \\
Proxmox VE & Free/Open Source & 0\,€ & 0\,€ \\
Metasploitable VMs & Free/Open Source & 0\,€ & 0\,€ \\
DVWA & Free/Open Source & 0\,€ & 0\,€ \\
Visual Studio Code & Free/Open Source & 0\,€ & 0\,€ \\
LaTeX (TeX Live) & Free/Open Source & 0\,€ & 0\,€ \\
Git + GitHub (Student) & Free/Student & 0\,€ & 0\,€ \\
Burp Suite Community & Free & 0\,€ & 0\,€ \\
Wireshark & Free/Open Source & 0\,€ & 0\,€ \\
\textbf{SUBTOTAL Software} & & & \textbf{0\,€} \\
\end{longtable}}

\subsubsection{Other Operational Costs}

{\footnotesize\begin{longtable}[]{@{}>{\raggedright\arraybackslash}p{0.45\linewidth}
  >{\raggedright\arraybackslash}p{0.45\linewidth}@{}}
\toprule\noalign{}
\textbf{Concept} & \textbf{Cost} \\
\midrule\noalign{}
\endhead
\bottomrule\noalign{}
\endlastfoot
Electricity (820\,h × 0.15\,kWh × 0.20\,€/kWh) & 25\,€ \\
Internet connectivity (9 months × 40\,€/month) & 360\,€ \\
Documentation printing and binding (3 copies) & 60\,€ \\
Contingency fund (5 & \\
\textbf{TOTAL Other Costs} & \textbf{1,}080\,€ \\
\end{longtable}}

\textbf{Note}: The total budget summary in Chapter (Table ) uses a rounded figure of 1,020\,€ for other costs, reflecting the exclusion of the contingency fund from the base estimate.
